\chapter{\texorpdfstring{$\beta$}{बीटा}-विलोपन}\label{ch:b1:elimination}

2-ब्रोमो-2-मेथिलप्रोपेन को कमरे के ताप पर एथेनॉल में सोडियम एथॉक्साइड के साथ
उपचारित कीजिए, तो मुख्य उत्पाद एक ईथर होता है; उसी मिश्रण को गरम कीजिए, तो
मुख्य उत्पाद एक ऐल्कीन, 2-मेथिलप्रोपीन, होता है, जिसमें कोई एथॉक्सी समूह
नहीं है। एथॉक्साइड आयन ने कार्बन पर आक्रमण करने वाले \omterm{def:b1:electronic-effects:nucleophile}{नाभिकरागी} की तरह नहीं, बल्कि
क्षारक की तरह कार्य किया है, जिसने ब्रोमाइड के निकलते समय बगल वाले कार्बन से
एक प्रोटॉन हटा दिया। प्रतिस्थापन और विलोपन के बीच यह प्रतिस्पर्धा
पूरे कार्बनिक संश्लेषण में व्याप्त है। यह अध्याय दोनों विलोपन
क्रियाविधियों, उनकी उल्लेखनीय ज्यामितीय अपेक्षा, यह पूर्वानुमान करने वाले नियम कि
कौन-सी ऐल्कीन बनेगी, और किसी अभिक्रिया को एक या दूसरी दिशा में मोड़ने के
ढंग का वर्णन करता है।

\begin{recall}
SN1 और SN2 (\cref{ch:b1:nucleophilic-substitution}); \omterm{def:b1:stereochemistry-in-depth:representations}{न्यूमैन प्रक्षेप},
साइक्लोहेक्सेन के \omterm{def:b1:stereochemistry-in-depth:isomers}{संरूपण} और $Z$/$E$ वर्णक
(\cref{ch:b1:stereochemistry-in-depth}); \omterm{def:b1:electronic-effects:intermediates}{कार्बधनायन} और क्षारक
(\cref{ch:b1:electronic-effects})।
\end{recall}

\section{E2 क्रियाविधि}

\begin{definition}[\texorpdfstring{$\beta$}{बीटा}-विलोपन, E2 और E1]\label{def:b1:elimination:elimination}
\emph{$\beta$-विलोपन}\index{बीटा-विलोपन@$\beta$-विलोपन} एक कार्बन ($\alpha$ कार्बन) से
\omterm{def:b1:electronic-effects:nucleophile}{निर्गामी समूह} Y को और पड़ोसी कार्बन (एक $\beta$ कार्बन) से एक प्रोटॉन को
हटाता है, जिससे उनके बीच द्वि-आबंध बनता है। \emph{E2 क्रियाविधि}\index{E2 क्रियाविधि} में
एक क्षारक प्रोटॉन लेता है जबकि Y निकलता है, एक ही \omterm{def:b1:elementary-steps:elementary-step}{मूल चरण} में; \emph{E1
क्रियाविधि}\index{E1 क्रियाविधि} में Y पहले निकलता है, जिससे एक \omterm{def:b1:electronic-effects:intermediates}{कार्बधनायन} बनता है, जो
फिर एक $\beta$ प्रोटॉन खो देता है।
\end{definition}

\begin{proposition}[E2 का दर नियम]\label{prop:b1:elimination:e2-law}
E2 के लिए $v = k[\mathrm{RY}][\mathrm{B}]$, जहाँ B क्षारक है।
\end{proposition}

\begin{proof}
\omterm{def:b1:nucleophilic-substitution:substitution}{क्रियाधार} और क्षारक को शामिल करने वाला एक द्विअणुक \omterm{def:b1:elementary-steps:elementary-step}{मूल चरण}: \omterm{def:b1:elementary-steps:elementary-step}{मूल चरणों} का
द्रव्य-अनुपाती क्रिया का नियम (\cref{ch:b1:elementary-steps})।
\end{proof}

\begin{definition}[प्रति- और सम-परिसमतलीय]\label{def:b1:elimination:periplanar}
पड़ोसी परमाणुओं पर स्थित दो आबंध \emph{प्रति-परिसमतलीय}\index{प्रति-परिसमतलीय}
होते हैं जब वे एक तल में विपरीत ओर हों (\omterm{def:b1:stereochemistry-in-depth:conformations}{मरोड़ कोण} $180^\circ$), और
\emph{सम-परिसमतलीय}\index{सम-परिसमतलीय} जब वे एक तल में
एक ही ओर हों ($0^\circ$)।
\end{definition}

\begin{omfigure}
\begin{tikzpicture}[font=\small, >={Stealth}]
  \pic (n) at (0,0) {omnewman={90}{270}};
  \node at (n-f1) {H}; \node at (n-f2) {\ce{CH3}}; \node at (n-f3) {H};
  \node at (n-b1) {Br}; \node at (n-b2) {\ce{CH3}}; \node at (n-b3) {H};
  \node (b) at (-2.6,2.0) {\ce{EtO-}};
  \draw[->, thick, omThm] (b) to[bend left=20] ($(n-f1)+(-0.15,0.1)$);
  \draw[->, thick, omThm] (0.12,0.8) to[bend left=40] (0.12,0.25);
  \draw[->, thick, omThm] (0.12,-0.6) to[bend right=35] ($(n-b1)+(0.25,0.15)$);
  \node[below] at (0,-2.0) {H और Br प्रति-परिसमतलीय};
  \draw[->, thick] (2.4,0) -- (3.6,0) node[midway, above] {E2};
  \node at (5.6,0) {\chemfig{C(-[:150]H_3C)(-[:210]H)=C(-[:30]H)(-[:-30]CH_3)}};
  \node at (5.6,-0.9) {\cip{E}-ब्यूट-2-ईन};
\end{tikzpicture}

{\small 2-ब्रोमोब्यूटेन की E2 अभिक्रिया, C3--C2 आबंध की दिशा में देखी गई (C3
सामने, Br वहन करने वाला C2 पीछे), उसके दो \omterm{def:b1:elimination:periplanar}{प्रति-परिसमतलीय}
\omterm{def:b1:stereochemistry-in-depth:isomers}{संरूपणों} में से एक में; दूसरा, जिसमें C3 का दूसरा हाइड्रोजन Br के प्रति प्रति-स्थिति में है,
\cip{Z}-ब्यूट-2-ईन देता है। तीन इलेक्ट्रॉन
युग्म एक साथ चलते हैं: क्षारक प्रोटॉन लेता है, C--H युग्म
$\pi$ आबंध बन जाता है, C--Br युग्म ब्रोमीन के साथ निकल जाता है। \omterm{def:b1:stereochemistry-in-depth:representations}{न्यूमैन प्रक्षेप} में
एक-दूसरे के सामने वाले समूह द्वि-आबंध की एक ही ओर पहुँचते हैं।}
\end{omfigure}

\begin{proposition}[E2 प्रति और त्रिविमविशिष्ट है]\label{prop:b1:elimination:anti}
E2 उस \omterm{def:b1:stereochemistry-in-depth:isomers}{संरूपण} से आगे बढ़ता है जिसमें C--H और C--Y आबंध
\omterm{def:b1:elimination:periplanar}{प्रति-परिसमतलीय} होते हैं। यह त्रिविमविशिष्ट है: \omterm{def:b1:stereochemistry-in-depth:enantiomers}{अप्रतिबिंबरूपी} \omterm{def:b1:nucleophilic-substitution:substitution}{क्रियाधार}
भिन्न ऐल्कीन \omterm{def:b1:stereochemistry-in-depth:isomers}{त्रिविम समावयवी} देते हैं, जो प्रति व्यवस्था द्वारा निर्धारित होते हैं।
\end{proposition}

\begin{proof}
\omterm{def:b1:elementary-steps:energy-profile}{संक्रमण अवस्था} में C--H \omterm{def:b1:lewis-resonance-vsepr:lewis}{आबंधी युग्म} को C--Y आबंध के प्रतिआबंधी क्षेत्र में
प्रवेश करना होता है जबकि दोनों कार्बन त्रिकोणीय बनते हैं; अतिव्यापन
सर्वोत्तम तब होता है जब दोनों आबंध समांतर हों, एक तल में। प्रति व्यवस्था
सांतरित है (निम्न ऊर्जा) और क्षारक को \omterm{def:b1:electronic-effects:nucleophile}{निर्गामी समूह} से दूर रखती है; सम व्यवस्था
ग्रसित है। H और Y प्रति-स्थिति में होने पर, हर कार्बन के दो अन्य समूह
निर्धारित हो जाते हैं: जो \omterm{def:b1:stereochemistry-in-depth:representations}{न्यूमैन प्रक्षेप} में एक ही ओर हैं, वे द्वि-आबंध पर
\textit{cis} पहुँचते हैं। किसी \omterm{def:b1:stereochemistry-in-depth:enantiomers}{अप्रतिबिंबरूपी} में एक कार्बन पर दो समूहों की अदला-बदली
होती है, और वह दूसरी ऐल्कीन देता है।
\end{proof}

\begin{proposition}[साइक्लोहेक्सेन पर E2]\label{prop:b1:elimination:cyclohexane-e2}
साइक्लोहेक्सेन वलय पर, E2 के लिए आवश्यक है कि \omterm{def:b1:electronic-effects:nucleophile}{निर्गामी समूह} और $\beta$
हाइड्रोजन दोनों \omterm{def:b1:stereochemistry-in-depth:conformations}{अक्षीय} हों (\textit{trans}-द्विअक्षीय); विषुवतीय स्थिति में रखा गया
\omterm{def:b1:electronic-effects:nucleophile}{निर्गामी समूह} तब तक अभिक्रिया नहीं कर सकता जब तक वलय पलट न जाए।
\end{proposition}

\begin{proof}
\omterm{def:b1:stereochemistry-in-depth:conformations}{कुर्सी} के सन्निकट कार्बनों पर दो \omterm{def:b1:stereochemistry-in-depth:conformations}{अक्षीय} आबंध \omterm{def:b1:elimination:periplanar}{प्रति-परिसमतलीय} होते हैं (एक
ऊपर, एक नीचे, अक्ष के समांतर); \omterm{def:b1:stereochemistry-in-depth:conformations}{विषुवतीय} आबंध केवल वलय के
C--C आबंधों के प्रति प्रति-स्थिति में होता है, किसी C--H आबंध के प्रति कभी नहीं।
\end{proof}

\begin{omfigure}
\begin{tikzpicture}[font=\small]
  \pic (a) at (0,0) {omchair};
  \draw[thick, omThm] (a-c1) -- (a-a1) node[above] {Cl};
  \draw[thick] (a-c1) -- (a-e1) node[right] {H};
  \draw[thick, omDef] (a-c2) -- (a-a2) node[below] {H};
  \draw[thick, omDef] (a-c6) -- (a-a6) node[below] {H};
  \node[align=left, anchor=west] at (3.0,0.2) {C1 पर Cl अक्षीय (ऊपर);\\दोनों पड़ोसियों पर अक्षीय H (नीचे):\\दो प्रति-परिसमतलीय H};
\end{tikzpicture}

{\small \omterm{def:b1:stereochemistry-in-depth:conformations}{कुर्सी} पर \textit{trans}-द्विअक्षीय व्यवस्था: \omterm{def:b1:stereochemistry-in-depth:conformations}{अक्षीय} क्लोरीन
दोनों पड़ोसी कार्बनों के \omterm{def:b1:stereochemistry-in-depth:conformations}{अक्षीय} हाइड्रोजनों (नीला) के प्रति
\omterm{def:b1:elimination:periplanar}{प्रति-परिसमतलीय} है।}
\end{omfigure}

\section{E1 क्रियाविधि}

\begin{proposition}[E1 का दर नियम]\label{prop:b1:elimination:e1-law}
E1 के लिए $v = k[\mathrm{RY}]$, दर आयनन की दर होती है।
\end{proposition}

\begin{proof}
जैसे SN1 के लिए (\cref{prop:b1:nucleophilic-substitution:sn1-law}):
\omterm{def:b1:electronic-effects:intermediates}{कार्बधनायन} एक मंद चरण में बनता है और तेज़ी से खपता है, यहाँ विलायक या किसी \omterm{def:b1:predominant-reaction:strong-acid}{दुर्बल
क्षारक} को प्रोटॉन देकर; \omterm{def:b1:elementary-steps:qssa}{स्थायी-अवस्था सन्निकटन}
$v = k_1[\mathrm{RY}]$ देता है।
\end{proof}

E1 अपना \omterm{def:b1:electronic-effects:intermediates}{कार्बधनायन} SN1 के साथ साझा करता है: \omterm{def:b1:intermolecular-forces-solvents:solvent-classes}{प्रोटिक विलायकों} में तृतीयक
\omterm{def:b1:nucleophilic-substitution:substitution}{क्रियाधार} दोनों देते हैं, और गरम करने से ऐल्कीन को वरीयता मिलती है। चूँकि
\omterm{def:b1:electronic-effects:intermediates}{कार्बधनायन} समतलीय है और प्रोटॉन बाद में निकलता है, E1 की कोई
ज्यामितीय अपेक्षा नहीं है और वह त्रिविमविशिष्ट नहीं है; वह मुख्यतः अधिक
स्थायी ऐल्कीन देता है, प्रायः \cip{E} वाली।

\begin{omfigure}
\setchemfig{atom sep=2.1em}
\schemestart
\chemfig{C(-[2]CH_3)(-[4]H_3C)(-[6]CH_3)-Br}
\arrow{->[मंद]}
\chemfig{C^{+}(-[2]CH_3)(-[:210]H_3C)(-[:-30]CH_3)}
\arrow{->[तीव्र][\ce{-H+}]}
\chemfig{C(-[2]CH_3)(-[:210]H_3C)=[:-30]CH_2}
\schemestop
\par\vspace{4pt}

{\small E1: 2-ब्रोमो-2-मेथिलप्रोपेन का \omterm{def:b1:electronic-effects:intermediates}{कार्बधनायन} तक आयनन, फिर
एक मेथिल समूह से विलायक को एक प्रोटॉन की हानि, जिससे
2-मेथिलप्रोपीन बनती है।}
\end{omfigure}

\section{स्थानवरणात्मकता: ज़ैत्सेव का नियम}

\begin{definition}[स्थानवरणात्मक और त्रिविमवरणात्मक अभिक्रियाएँ]\label{def:b1:elimination:selectivity}
जो अभिक्रिया कई \omterm{def:b1:stereochemistry-in-depth:isomers}{संरचनात्मक समावयवी} दे सकती है पर मुख्यतः एक बनाती है,
वह \emph{स्थानवरणात्मक अभिक्रिया}\index{स्थानवरणात्मक अभिक्रिया} है;
जो कई \omterm{def:b1:stereochemistry-in-depth:isomers}{त्रिविम समावयवी} दे सकती है पर मुख्यतः एक बनाती है, वह
\emph{त्रिविमवरणात्मक अभिक्रिया}\index{त्रिविमवरणात्मक अभिक्रिया} है।
\end{definition}

\begin{proposition}[ज़ैत्सेव का नियम]\label{prop:b1:elimination:zaitsev}
जब कई $\beta$ हाइड्रोजन हटाए जा सकते हों, तो मुख्य ऐल्कीन
प्रायः सबसे अधिक प्रतिस्थापित (सबसे स्थायी) होती है, और \cip{Z} के बजाय \cip{E} समावयवी
(\emph{ज़ैत्सेव का नियम}\index{ज़ैत्सेव का नियम})।
भारी-भरकम क्षारक, और ज्यामितीय बाध्यताएँ जैसे trans-द्विअक्षीय
अपेक्षा, इसे निष्प्रभावी कर सकते हैं।
\end{proposition}

\begin{proof}
द्वि-आबंध के कार्बनों पर ऐल्किल समूह उसे स्थायी बनाते हैं (जैसे वे
\omterm{def:b1:electronic-effects:intermediates}{कार्बधनायनों} को स्थायी बनाते हैं), और \cip{E} ऐल्कीन \textit{cis} समूहों की
भीड़ से बचती हैं। E2 और E1 की \omterm{def:b1:elementary-steps:energy-profile}{संक्रमण अवस्थाओं} में आंशिक
द्वि-आबंध लक्षण होता है, इसलिए अधिक स्थायी ऐल्कीन तेज़ी से बनती है।
भारी-भरकम क्षारक कम बाधित हाइड्रोजनों तक अधिक आसानी से पहुँचता है; और जो
हाइड्रोजन \omterm{def:b1:electronic-effects:nucleophile}{निर्गामी समूह} के प्रति \omterm{def:b1:elimination:periplanar}{प्रति-परिसमतलीय} नहीं हो सकता, उसे
E2 बिल्कुल नहीं हटाता, चाहे वह जो ऐल्कीन देता उसका स्थायित्व कुछ भी हो।
\end{proof}

\begin{method}[ऐल्कीन का पूर्वानुमान]\label{met:b1:elimination:predict-alkene}
\begin{enumerate}
  \item हाइड्रोजन वहन करने वाले $\beta$ कार्बनों की सूची बनाइए।
  \item E2 के लिए केवल वे हाइड्रोजन रखिए जो Y के प्रति \omterm{def:b1:elimination:periplanar}{प्रति-परिसमतलीय} हो सकते हैं (वलय
        पर: trans-द्विअक्षीय, ऐसी \omterm{def:b1:stereochemistry-in-depth:conformations}{कुर्सी} में जो वास्तव में बन सकती हो)।
  \item उनमें से उसे वरीयता दीजिए जो सबसे अधिक प्रतिस्थापित ऐल्कीन देता है
        (ज़ैत्सेव), जब तक क्षारक भारी-भरकम न हो।
  \item हर एक के लिए, \omterm{def:b1:stereochemistry-in-depth:conformations}{प्रति संरूपण} में \omterm{def:b1:stereochemistry-in-depth:representations}{न्यूमैन प्रक्षेप} खींचकर
        \cip{E}/\cip{Z} ज्यामिति पढ़िए।
\end{enumerate}
\end{method}

\section{प्रतिस्थापन या विलोपन?}

\begin{proposition}[प्रतिस्थापन बनाम विलोपन]\label{prop:b1:elimination:sn-vs-e}
विलोपन को प्रबल, भारी-भरकम क्षारकों, उच्च ताप और
तृतीयक या भीड़ वाले \omterm{def:b1:nucleophilic-substitution:substitution}{क्रियाधारों} से वरीयता मिलती है; प्रतिस्थापन को ऐसे अच्छे \omterm{def:b1:electronic-effects:nucleophile}{नाभिकरागियों} से जो
\omterm{def:b1:predominant-reaction:strong-acid}{दुर्बल क्षारक} हों, निम्न ताप से और प्राथमिक \omterm{def:b1:nucleophilic-substitution:substitution}{क्रियाधारों} से।
\end{proposition}

\begin{proof}
क्षारक की प्रबलता हाइड्रोजन पर आक्रमण के अनुकूल है; भारीपन
भीड़ वाले कार्बन पर आक्रमण में बाधा डालता है, पर खुले हाइड्रोजन पर नहीं। \omterm{def:b1:electronic-effects:carbon-class}{तृतीयक कार्बन}
SN2 के लिए अगम्य है, और उसका \omterm{def:b1:electronic-effects:intermediates}{कार्बधनायन} आसानी से प्रोटॉन खो देता है। विलोपन
एक अणु को दो या तीन में बदल देता है (ऐल्कीन, \omterm{def:b1:electronic-effects:nucleophile}{निर्गामी समूह}, प्रोटॉनित
क्षारक): ताप बढ़ने के साथ इसे वरीयता मिलती है, जिसका औचित्य
द्वितीय वर्ष के खंड में दिया गया है।
\end{proof}

\begin{omfigure}
\begin{tabular}{llll}
\hline
\omterm{def:b1:nucleophilic-substitution:substitution}{क्रियाधार} & \omterm{def:b1:predominant-reaction:strong-acid}{प्रबल क्षारक}, छोटा & \omterm{def:b1:predominant-reaction:strong-acid}{प्रबल क्षारक}, भारी-भरकम & \omterm{def:b1:predominant-reaction:strong-acid}{दुर्बल क्षारक}, अच्छा \omterm{def:b1:electronic-effects:nucleophile}{नाभिकरागी} \\
\hline
प्राथमिक & SN2 (कुछ E2) & E2 & SN2 \\
द्वितीयक & E2 (और SN2) & E2 & SN2 (अप्रोटिक), SN1/E1 (प्रोटिक) \\
तृतीयक & E2 & E2 & SN1/E1 (प्रोटिक) \\
\hline
\end{tabular}

\smallskip
{\small एक प्रारंभिक दिशा-निर्देश: हर स्थिति में मुख्य अभिक्रिया। गरम करना
हर स्तंभ में विलोपन के अनुकूल है।}
\end{omfigure}

\section{अभ्यास}

\begin{exercise}[$\star$]\label{exo:b1:elimination:1}
2-ब्रोमोब्यूटेन, 2-ब्रोमो-2-मेथिलब्यूटेन और ब्रोमोसाइक्लोहेक्सेन से
HBr के विलोपन द्वारा बन सकने वाली ऐल्कीन बताइए, और वह भी जिसे \omterm{prop:b1:elimination:zaitsev}{ज़ैत्सेव का
नियम} मुख्य उत्पाद के रूप में बताता है।
\end{exercise}

\begin{exercise}[$\star$]\label{exo:b1:elimination:2}
E1 और E2 के \omterm{def:b1:rate-laws:order}{दर नियम} लिखिए। एथॉक्साइड युक्त एथेनॉल में 2-ब्रोमो-2-मेथिलप्रोपेन के लिए
आप प्रायोगिक रूप से उन्हें कैसे अलग पहचानेंगे?
\end{exercise}

\begin{exercise}[$\star$]\label{exo:b1:elimination:3}
C2--C3 की दिशा में 2-ब्रोमोब्यूटेन के तीनों \omterm{def:b1:stereochemistry-in-depth:conformations}{सांतरित संरूपणों} के \omterm{def:b1:stereochemistry-in-depth:representations}{न्यूमैन
प्रक्षेप} खींचिए, और बताइए कि किसमें E2 हो सकता है।
\end{exercise}

\begin{exercise}[$\star$]\label{exo:b1:elimination:4}
मुख्य अभिक्रिया (SN1, SN2, E1, E2) का पूर्वानुमान कीजिए: \qty{25}{\celsius} पर जल में
सोडियम हाइड्रॉक्साइड के साथ ब्रोमोएथेन; पोटैशियम टर्ट-ब्यूटॉक्साइड के साथ
2-ब्रोमो-2-मेथिलप्रोपेन; प्रोपेनोन में सोडियम आयोडाइड के साथ 2-ब्रोमोप्रोपेन;
गरम एथेनॉल में 2-ब्रोमो-2-मेथिलप्रोपेन।
\end{exercise}

\begin{exercise}[$\star\star$]\label{exo:b1:elimination:5}
\omterm{def:b1:stereochemistry-in-depth:representations}{न्यूमैन प्रक्षेपों} से दिखाइए कि 2-ब्रोमो-3-मेथिलपेंटेन के वे दो \omterm{def:b1:stereochemistry-in-depth:enantiomers}{अप्रतिबिंबरूपी} जो
C3 हाइड्रोजन खो सकते हैं, HBr के E2 विलोपन से
3-मेथिलपेंट-2-ईन के दोनों \omterm{def:b1:stereochemistry-in-depth:isomers}{त्रिविम समावयवी} देते हैं।
\end{exercise}

\begin{exercise}[$\star\star$]\label{exo:b1:elimination:6}
गरम सांद्र सल्फ़्यूरिक अम्ल में 2-मेथिलब्यूटेन-2-ऑल के निर्जलीकरण (E1) की क्रियाविधि
लिखिए और मुख्य ऐल्कीन बताइए।
\end{exercise}

\begin{exercise}[$\star\star$]\label{exo:b1:elimination:7}
समझाइए कि \textit{trans}-1-ब्रोमो-4-टर्ट-ब्यूटिलसाइक्लोहेक्सेन E2 से बहुत
धीरे अभिक्रिया क्यों करता है जबकि उसका \textit{cis} समावयवी तेज़ी से। (टर्ट-ब्यूटिल
समूह \omterm{def:b1:stereochemistry-in-depth:conformations}{विषुवतीय} रहता है।)
\end{exercise}

\begin{exercise}[$\star\star$]\label{exo:b1:elimination:8}
सोडियम एथॉक्साइड के साथ 2-ब्रोमो-2-मेथिलब्यूटेन मुख्यतः
2-मेथिलब्यूट-2-ईन देता है; पोटैशियम टर्ट-ब्यूटॉक्साइड के साथ, मुख्यतः
2-मेथिलब्यूट-1-ईन। व्याख्या कीजिए।
\end{exercise}

\begin{exercise}[$\star\star$]\label{exo:b1:elimination:9}
ब्रोमोमेथेन के साथ विलोपन कभी क्यों नहीं दिखता? और
1-ब्रोमो-2,2-डाइमेथिलप्रोपेन के साथ E2 द्वारा?
\end{exercise}

\begin{exercise}[$\star\star\star$]\label{exo:b1:elimination:10}
\qty{25}{\celsius} पर एथेनॉल में 2-ब्रोमो-2-मेथिलप्रोपेन ईथर (SN1) और
ऐल्कीन (E1) का ऐसा मिश्रण देता है जिसके अनुपात \omterm{def:b1:nucleophilic-substitution:substitution}{क्रियाधार} की
सांद्रता पर निर्भर नहीं करते। दिखाइए कि दोनों उत्पाद
एक ही \omterm{def:b1:electronic-effects:intermediates}{कार्बधनायन} से आते हैं और उनका अनुपात दो \omterm{def:b1:rate-laws:order}{दर
स्थिरांकों} का अनुपात है।
\end{exercise}

\begin{exercise}[$\star\star\star$]\label{exo:b1:elimination:11}
1,2-डाइब्रोमो-1,2-डाइफ़ेनिलएथेन का एक \omterm{def:b1:stereochemistry-in-depth:enantiomers}{अप्रतिबिंबरूपी} (मेसो वाला) एथॉक्साइड के साथ E2 द्वारा
1-ब्रोमो-1,2-डाइफ़ेनिलएथीन का एक अकेला \omterm{def:b1:stereochemistry-in-depth:isomers}{त्रिविम समावयवी} देता है।
\omterm{def:b1:stereochemistry-in-depth:representations}{न्यूमैन प्रक्षेप} की सहायता से ज्ञात कीजिए कि कौन-सा।
\end{exercise}

\begin{exercise}[$\star\star\star$]\label{exo:b1:elimination:12}
\cref{ch:b1:stereochemistry-in-depth} की ऊर्जाओं का उपयोग करके समझाइए कि
जब \omterm{def:b1:electronic-effects:nucleophile}{निर्गामी समूह} को भीड़ वाली \omterm{def:b1:stereochemistry-in-depth:conformations}{कुर्सी} में \omterm{def:b1:stereochemistry-in-depth:conformations}{अक्षीय} बनना पड़ता है, तो साइक्लोहेक्सेन पर
E2 अभिक्रिया सैकड़ों गुना धीमी क्यों हो सकती है। दर को
किसी संरूपी के अंश और एक \omterm{def:b1:rate-laws:order}{दर स्थिरांक} के गुणनफल के रूप में व्यक्त कीजिए।
\end{exercise}

\section{समस्या: मेंथिल और निओमेंथिल क्लोराइड}

\begin{problem}[{सप्ताहांत समस्या --- दो अप्रतिबिंबरूपी
क्लोराइडों की कुर्सियाँ, हर एक में उपलब्ध प्रति-परिसमतलीय हाइड्रोजन, बनने वाली
ऐल्कीन, और उनकी E2 दरों का अनुपात}]
\label{pb:b1:elimination:1}
मेंथिल क्लोराइड और निओमेंथिल क्लोराइड मेंथॉल
और निओमेंथॉल (\cref{ch:b1:stereochemistry-in-depth}) के क्लोराइड हैं: वलय पर C1
Cl वहन करता है, C2 आइसोप्रोपिल समूह, C5 मेथिल समूह। मेंथिल
क्लोराइड में तीनों समूह \omterm{def:b1:stereochemistry-in-depth:conformations}{विषुवतीय} हो सकते हैं; निओमेंथिल क्लोराइड में,
जब दोनों ऐल्किल समूह \omterm{def:b1:stereochemistry-in-depth:conformations}{विषुवतीय} हों तो Cl \omterm{def:b1:stereochemistry-in-depth:conformations}{अक्षीय} होता है। दोनों को
एथेनॉल में सोडियम एथॉक्साइड से उपचारित किया जाता है, एक E2 अभिक्रिया। समस्या के आँकड़े:
निओमेंथिल क्लोराइड का \qty{95}{\percent} भाग Cl \omterm{def:b1:stereochemistry-in-depth:conformations}{अक्षीय} वाली \omterm{def:b1:stereochemistry-in-depth:conformations}{कुर्सी} में है;
मेंथिल क्लोराइड के लिए, Cl \omterm{def:b1:stereochemistry-in-depth:conformations}{अक्षीय} वाली \omterm{def:b1:stereochemistry-in-depth:conformations}{कुर्सी} (तीनों समूह \omterm{def:b1:stereochemistry-in-depth:conformations}{अक्षीय})
का अंश \num{5.0e-3} है; मान लीजिए कि किसी क्रियाशील \omterm{def:b1:stereochemistry-in-depth:conformations}{कुर्सी} का हर
\omterm{def:b1:elimination:periplanar}{प्रति-परिसमतलीय} हाइड्रोजन समान \omterm{def:b1:rate-laws:order}{दर स्थिरांक} से अभिक्रिया करता है। निओमेंथिल क्लोराइड
\qty{78}{\percent} त्रिप्रतिस्थापित ऐल्कीन देता है।

\medskip
\noindent\textbf{भाग I --- \omterm{def:b1:stereochemistry-in-depth:conformations}{कुर्सियाँ}।}
\begin{enumerate}
  \item मेंथिल क्लोराइड की वरीय \omterm{def:b1:stereochemistry-in-depth:conformations}{कुर्सी} खींचिए।
  \item निओमेंथिल क्लोराइड की वरीय \omterm{def:b1:stereochemistry-in-depth:conformations}{कुर्सी} खींचिए।
  \item क्या दोनों क्लोराइड \omterm{def:b1:stereochemistry-in-depth:enantiomers}{प्रतिबिंबरूपी} हैं या \omterm{def:b1:stereochemistry-in-depth:enantiomers}{अप्रतिबिंबरूपी}?
  \item वलय पर E2 की ज्यामितीय अपेक्षा बताइए।
  \item मेंथिल क्लोराइड किस \omterm{def:b1:stereochemistry-in-depth:conformations}{कुर्सी} में अभिक्रिया कर सकता है? उसे खींचिए।
  \item वह \omterm{def:b1:stereochemistry-in-depth:conformations}{कुर्सी} दुर्लभ क्यों है? \omterm{def:b1:stereochemistry-in-depth:conformations}{अक्षीय} समूहों की ऊर्जा-कीमत का उपयोग कीजिए।
\end{enumerate}

\medskip
\noindent\textbf{भाग II --- प्रति हाइड्रोजन।}
\begin{enumerate}[resume]
  \item क्लोराइड के $\beta$ कार्बनों (C2 और C6) और उनके
        हाइड्रोजनों की सूची बनाइए।
  \item निओमेंथिल क्लोराइड की क्रियाशील \omterm{def:b1:stereochemistry-in-depth:conformations}{कुर्सी} में कौन-से हाइड्रोजन
        Cl के प्रति \omterm{def:b1:elimination:periplanar}{प्रति-परिसमतलीय} हैं?
  \item मेंथिल क्लोराइड की क्रियाशील \omterm{def:b1:stereochemistry-in-depth:conformations}{कुर्सी} में, क्या C2 का हाइड्रोजन
        \omterm{def:b1:stereochemistry-in-depth:conformations}{अक्षीय} है? क्या C6 का कोई हाइड्रोजन \omterm{def:b1:stereochemistry-in-depth:conformations}{अक्षीय} है?
  \item हर क्लोराइड कितने \omterm{def:b1:elimination:periplanar}{प्रति-परिसमतलीय} हाइड्रोजन प्रस्तुत करता है?
  \item मेंथिल क्लोराइड की क्रियाशील \omterm{def:b1:stereochemistry-in-depth:conformations}{कुर्सी} के लिए C1--C6 की दिशा में
        \omterm{def:b1:stereochemistry-in-depth:representations}{न्यूमैन प्रक्षेप} खींचिए।
  \item \omterm{def:b1:elimination:periplanar}{सम-परिसमतलीय} विलोपन अभिक्रिया को क्यों नहीं बचा
        सकता?
\end{enumerate}

\medskip
\noindent\textbf{भाग III --- उत्पाद।}
\begin{enumerate}[resume]
  \item C2 हाइड्रोजन हटाकर बनने वाली ऐल्कीन का और
        C6 के एक हाइड्रोजन को हटाकर बनने वाली ऐल्कीन का नाम बताइए।
  \item \omterm{prop:b1:elimination:zaitsev}{ज़ैत्सेव का नियम} किसे वरीयता देता है, और क्यों?
  \item निओमेंथिल क्लोराइड क्या देता है? क्या 78/22 अनुपात
        \omterm{prop:b1:elimination:zaitsev}{ज़ैत्सेव के नियम} से संगत है?
  \item मेंथिल क्लोराइड क्या देता है? क्या यह
        \omterm{prop:b1:elimination:zaitsev}{ज़ैत्सेव के नियम} से संगत है?
  \item कौन-सी अभिक्रिया त्रिविमविशिष्ट है, कौन-सी स्थानवरणात्मक, या दोनों?
  \item यहाँ एथॉक्साइड के साथ प्रतिस्थापन (SN2) गौण क्यों है?
\end{enumerate}

\medskip
\noindent\textbf{भाग IV --- दरें।}
\begin{enumerate}[resume]
  \item हर अभिक्रिया की दर को क्रियाशील \omterm{def:b1:stereochemistry-in-depth:conformations}{कुर्सी} के अंश और
        \omterm{def:b1:elimination:periplanar}{प्रति-परिसमतलीय} हाइड्रोजनों की संख्या के फलन के रूप में लिखिए।
  \item निओमेंथिल क्लोराइड की आपेक्षिक दर परिकलित कीजिए।
  \item मेंथिल क्लोराइड की आपेक्षिक दर परिकलित कीजिए।
  \item दरों का अनुपात $k(\text{निओमेंथिल})/k(\text{मेंथिल})$ परिकलित कीजिए।
\end{enumerate}
\end{problem}
